\documentclass[12pt,letterpaper]{article}
\usepackage[OT1]{fontenc}
\renewcommand{\rmdefault}{cmr}
\renewcommand{\sfdefault}{cmss}
\renewcommand{\ttdefault}{cmtt}
\usepackage{amsmath,amssymb}
\usepackage[left=1.25in,right=1.25in,top=1.05in,bottom=1.38in,footskip=0.40in]{geometry}
\usepackage{enumitem,graphicx,needspace}
\setlength{\parindent}{0pt}
\setlength{\parskip}{0.45em}
\setlist{topsep=0.35em,itemsep=0.3em,parsep=0pt}
\setlength{\emergencystretch}{2em}
\pagestyle{plain}
\newcommand{\problemheading}[2][5\baselineskip]{\Needspace{#1}\par\noindent\hspace*{-0.625in}\textbf{#2}\par\nobreak\vspace{0.3em}}
\newcommand{\examspace}[1]{\par\vspace*{#1cm}}
\begin{document}
\thispagestyle{empty}
\begin{center}
\textbf{COMP 2711}\\
\textbf{Discrete Math Tools for Computer Science}\\
{\large\bfseries Mock Midterm Examination 10}\\[1em]
\textbf{Practice material -- not an official examination}\\[0.7em]
Time allowed: 120 minutes \hfill Total: 100 points
\end{center}
\vspace{1.1cm}
\begin{center}\fbox{\fbox{\begin{minipage}{0.93\textwidth}
Name: \hrulefill\quad Student ID: \hrulefill\\[0.6em]
Email: \hrulefill\quad Lecture: \hrulefill
\end{minipage}}}\end{center}
\vspace{0.6cm}
\textbf{Instructions}
\begin{itemize}[leftmargin=1.3em]
\item This closed-book practice exam consists of 10 problems. You may use one A4 single-sided handwritten cheat sheet.
\item Answer all problems in the spaces provided. You may also use the backs of pages and the draft pages. Each problem starts on a new page; some answers require only a few lines.
\item Give a proof, justification, or counterexample when requested. State any additional notation you introduce.
\item Answers may be left in exact form using binomial coefficients, factorials, and $C(n,k)$ or $P(n,k)$. For example, $\binom{5}{3}+\binom{4}{2}$ is acceptable. Avoid nonstandard notation such as $nCk$ or $nPk$.
\item The separate solutions document contains derivations and marking notes. Attempt the exam before consulting it.
\end{itemize}
\vspace{1.0cm}
\begin{center}
\Large
\setlength{\tabcolsep}{2pt}
\begin{tabular*}{\textwidth}{@{\extracolsep{\fill}}|l|*{11}{c|}}\hline
Problems & 1 & 2 & 3 & 4 & 5 & 6 & 7 & 8 & 9 & 10 & Total\\\hline
Points & 10 & 13 & 12 & 16 & 10 & 9 & 7 & 6 & 7 & 10 & 100\\\hline
Score & & & & & & & & & & & \\[0.7em]\hline
\end{tabular*}
\end{center}
\vspace{0.8cm}

\clearpage
\vspace*{0.21in}
\problemheading{Problem 1: Propositional logic and proofs [10 points]}
\begin{list}{}{\setlength{\leftmargin}{0.25in}\setlength{\rightmargin}{0.25in}}\item[]
\begin{description}[font=\normalfont,leftmargin=1.2em,labelsep=0.4em]
\item[(a)] \textbf{[3 points]} Classify with a logical-law justification: (i) $(p\to q)\lor(q\to p)$; (ii) $(p\leftrightarrow q)\land(p\oplus q)$; (iii) $(p\to q)\land\neg q$. Each is worth 1 point.
\examspace{4}
\Needspace{\dimexpr 3cm+4\baselineskip\relax}
\item[(b)] \textbf{[2 points]} Is the liar sentence \emph{This sentence is false} a proposition under the classical definition used in the lecture (1 point)? If $p\to q$ is given, does that logically assert \emph{$q$ only if $p$} (1 point)? Explain briefly.
\examspace{3}
\Needspace{\dimexpr 6cm+4\baselineskip\relax}
\item[(c)] \textbf{[5 points]} Give direct proofs: (i) if $u,v$ are rational and $v\ne0$, then $u/v$ is rational (3 points); (ii) if an integer $n$ is odd, then $n(n+2)$ is odd (2 points).
\examspace{6}
\end{description}
\end{list}
\clearpage
\vspace*{0.21in}
\problemheading{Problem 2: Predicate logic [13 points]}
\begin{list}{}{\setlength{\leftmargin}{0.25in}\setlength{\rightmargin}{0.25in}}\item[]
\begin{description}[font=\normalfont,leftmargin=1.2em,labelsep=0.4em]
\item[(a)] \textbf{[4 points]} In $\forall xP(x)\lor Q(x)$, use quantifier precedence to insert parentheses, identify free occurrences and decide whether it is a proposition (2 points). Translate \emph{some real is at least every real} and give its truth value with justification (2 points).
\examspace{4}
\Needspace{\dimexpr 5cm+4\baselineskip\relax}
\item[(b)] \textbf{[4 points]} Let $A$ be a proposition with no free $x$, and first assume the domain is nonempty. Prove $\forall x(P(x)\to A)\equiv((\exists xP(x))\to A)$ (2 points). Now let the domain be empty: give a truth value of $A$ showing that $\forall x(A\land P(x))\equiv A\land\forall xP(x)$ need not hold (2 points).
\examspace{5}
\Needspace{\dimexpr 6cm+4\baselineskip\relax}
\item[(c)] \textbf{[5 points]} From $\forall x(P(x)\to Q(x))$, $\forall x(R(x)\to\neg Q(x))$, and $\exists x(R(x)\land S(x))$, prove $\exists x(S(x)\land\neg P(x))$ using inference rules.
\examspace{6}
\end{description}
\end{list}
\clearpage
\vspace*{0.21in}
\problemheading{Problem 3: Sets and functions [12 points]}
\begin{list}{}{\setlength{\leftmargin}{0.25in}\setlength{\rightmargin}{0.25in}}\item[]
Use $\mathbb N=\{0,1,2,\ldots\}$.
\begin{description}[font=\normalfont,leftmargin=1.2em,labelsep=0.4em]
\item[(a)] \textbf{[2 points]} Within $U=\{0,1,2,3,4,5\}$, let $A=\{0,2,4\}$ and $B=\{2,3,4\}$. Find $A\setminus B$, $U\setminus(A\cup B)$, and $|A\cup B|$, verifying the last value by the two-set identity.
\examspace{3}
\Needspace{\dimexpr 5cm+4\baselineskip\relax}
\item[(b)] \textbf{[3 points]} Construct an explicit bijection $h:\mathbb N\to\mathbb Z$ by assigning nonnegative integers to even inputs and negative integers to odd inputs. Prove injectivity and surjectivity.
\examspace{5}
\Needspace{\dimexpr 1cm+4\baselineskip\relax}
\item[(c)] \textbf{[7 points]} Write T or F for each statement. No justification is required. Each item is worth 1 point.\begin{enumerate}[label=(\roman*)]\item $\mathbb N=\{0,1,2,\ldots\}$ is finite.\item The nonnegative even integers are countably infinite.\item The set of binary sequences with exactly six ones is countably infinite.\item $\mathbb Q\cap(0,1)$ is countably infinite.\item Every proper subset of a finite set has smaller cardinality.\item An uncountable set can be listed by the natural numbers with no omissions.\item The set of all finite computer programs is uncountable.\end{enumerate}
\examspace{1}
\end{description}
\end{list}
\clearpage
\vspace*{0.21in}
\problemheading{Problem 4: Short counting calculations [16 points]}
\begin{list}{}{\setlength{\leftmargin}{0.25in}\setlength{\rightmargin}{0.25in}}\item[]
Give exact answers. Factorials, powers, and binomial coefficients may be left unevaluated. No explanation is required unless requested.
\begin{description}[font=\normalfont,leftmargin=1.2em,labelsep=0.4em]
\item[(a)] \textbf{[2 points]} Evaluate $\sum_{k=0}^{6}\binom6k^2$ using a binomial identity.
\examspace{1.5}
\Needspace{\dimexpr 1.5cm+4\baselineskip\relax}
\item[(b)] \textbf{[2 points]} Eight distinct people sit around a circle with three specified people consecutive in any internal order. Rotations are equivalent; reflections are different. Count the seatings.
\examspace{1.5}
\Needspace{\dimexpr 1.5cm+4\baselineskip\relax}
\item[(c)] \textbf{[2 points]} Count strings using exactly two $A$s, two $B$s and two $C$s, with the two $A$s not adjacent. Equal symbols are indistinguishable.
\examspace{1.5}
\Needspace{\dimexpr 1.5cm+4\baselineskip\relax}
\item[(d)] \textbf{[2 points]} Choose an unordered six-person committee from ten distinct people, with $A$ and $B$ not both selected. Count the committees.
\examspace{1.5}
\Needspace{\dimexpr 1.5cm+4\baselineskip\relax}
\item[(e)] \textbf{[2 points]} Choose six items from three unlimited types $A,B,C$, with an even number of type $A$. Items of a type are identical and order does not matter. Count the choices.
\examspace{1.5}
\Needspace{\dimexpr 1.5cm+4\baselineskip\relax}
\item[(f)] \textbf{[2 points]} Find the constant term in $(x-2x^{-1})^8$.
\examspace{1.5}
\Needspace{\dimexpr 1.5cm+4\baselineskip\relax}
\item[(g)] \textbf{[2 points]} Count injective functions $f:\{1,2,3,4\}\to\{1,2,3,4\}$ with $f(2)=2$ and $f(1)\ne1$.
\examspace{1.5}
\Needspace{\dimexpr 1.5cm+4\baselineskip\relax}
\item[(h)] \textbf{[2 points]} Count right/up unit-step lattice paths from $(0,0)$ to $(4,3)$ that do not begin with two right steps.
\examspace{1.5}
\end{description}
\end{list}
\clearpage
\vspace*{0.21in}
\problemheading{Problem 5: Combinatorial proof [10 points]}
\begin{list}{}{\setlength{\leftmargin}{0.25in}\setlength{\rightmargin}{0.25in}}\item[]
\begin{description}[font=\normalfont,leftmargin=1.2em,labelsep=0.4em]
\item[] \textbf{[10 points]} For $n\ge1$, use a complement bijection to prove\[\sum_{k=0}^{n-1}\binom{2n}{k}=\frac{2^{2n}-\binom{2n}{n}}{2}.\]Explain why the subsets of size $n$ must be treated separately.
\examspace{7}
\end{description}
\end{list}
\clearpage
\vspace*{0.21in}
\problemheading{Problem 6: Inclusion-exclusion [9 points]}
\begin{list}{}{\setlength{\leftmargin}{0.25in}\setlength{\rightmargin}{0.25in}}\item[]
\begin{description}[font=\normalfont,leftmargin=1.2em,labelsep=0.4em]
\item[] \textbf{[9 points]} Use inclusion--exclusion to count integers in $\{1,\ldots,180\}$ divisible by exactly two of $2,3,5$. Explain how your expression removes integers divisible by all three.
\examspace{7}
\end{description}
\end{list}
\clearpage
\vspace*{0.21in}
\problemheading{Problem 7: Pigeonhole principle [7 points]}
\begin{list}{}{\setlength{\leftmargin}{0.25in}\setlength{\rightmargin}{0.25in}}\item[]
\begin{description}[font=\normalfont,leftmargin=1.2em,labelsep=0.4em]
\item[] \textbf{[7 points]} Prove that any seventeen distinct bit strings of length eight include three with the same number of ones. Show that seventeen is the smallest number with this guarantee.
\examspace{7}
\end{description}
\end{list}
\clearpage
\vspace*{0.21in}
\problemheading{Problem 8: Joint masses and event rules [6 points]}
\begin{list}{}{\setlength{\leftmargin}{0.25in}\setlength{\rightmargin}{0.25in}}\item[]
An outcome is an ordered bit pair. Its probabilities are $P(00)=1/2$, $P(01)=1/4$, $P(10)=1/8$, $P(11)=1/8$. Let $A$ mean first bit one and $B$ mean second bit one.
\begin{description}[font=\normalfont,leftmargin=1.2em,labelsep=0.4em]
\item[(a)] \textbf{[1 points]} Find $P(A)$ and $P(B)$.
\examspace{1.2}
\Needspace{\dimexpr 1.2cm+4\baselineskip\relax}
\item[(b)] \textbf{[1 points]} Find $P(A\cap B)$.
\examspace{1.2}
\Needspace{\dimexpr 1.2cm+4\baselineskip\relax}
\item[(c)] \textbf{[1 points]} Find $P(A\cup B)$.
\examspace{1.2}
\Needspace{\dimexpr 1.2cm+4\baselineskip\relax}
\item[(d)] \textbf{[1 points]} Find $P(A^c\cap B^c)$.
\examspace{1.2}
\Needspace{\dimexpr 2.0cm+4\baselineskip\relax}
\item[(e)] \textbf{[2 points]} Justify the two-event union rule by a disjoint-set decomposition, and test independence of $A,B$.
\examspace{2.0}
\end{description}
\end{list}
\clearpage
\vspace*{0.21in}
\problemheading{Problem 9: Restricted bit strings [7 points]}
\begin{list}{}{\setlength{\leftmargin}{0.25in}\setlength{\rightmargin}{0.25in}}\item[]
Choose uniformly among the sixteen bit strings of length four. Let $E$ mean first bit one, $F$ mean exactly two ones, and $G$ mean that the two ones are consecutive (so $G\subseteq F$).
\begin{description}[font=\normalfont,leftmargin=1.2em,labelsep=0.4em]
\item[(a)] \textbf{[3 points]} Find $P(E\mid F)$ and $P(G\mid F)$ using the restricted sample space.
\examspace{3.5}
\Needspace{\dimexpr 4.5cm+4\baselineskip\relax}
\item[(b)] \textbf{[4 points]} Find $P(F\mid E)$ and use it to test independence of $E,F$. State the positivity requirement for this conditional test.
\examspace{4.5}
\end{description}
\end{list}
\clearpage
\vspace*{0.21in}
\problemheading{Problem 10: Bernoulli sums and variance proofs [10 points]}
\begin{list}{}{\setlength{\leftmargin}{0.25in}\setlength{\rightmargin}{0.25in}}\item[]
Let $I_1,\ldots,I_n$ be independent Bernoulli indicators with success probability $p$, where $n\ge1$ and $0<p<1$. Put $S=\sum_i I_i$ and $M=S/n$.
\begin{description}[font=\normalfont,leftmargin=1.2em,labelsep=0.4em]
\item[(a)] \textbf{[2 points]} Compute $E(I_i)$ and $V(I_i)$ from its two possible values.
\examspace{2.0}
\Needspace{\dimexpr 2.0cm+4\baselineskip\relax}
\item[(b)] \textbf{[2 points]} Find $E(S)$ and $V(S)$, identifying where independence is used.
\examspace{2.0}
\Needspace{\dimexpr 2.0cm+4\baselineskip\relax}
\item[(c)] \textbf{[2 points]} Find $E(M)$ and $V(M)$.
\examspace{2.0}
\Needspace{\dimexpr 4.5cm+4\baselineskip\relax}
\item[(d)] \textbf{[4 points]} Derive $V(X)=E(X^2)-E(X)^2$ for a finite random variable. Then prove $V(X+Y)=V(X)+V(Y)$ when $X,Y$ are independent finite variables.
\examspace{4.5}
\end{description}
\end{list}
\clearpage\textbf{Draft Paper}\par\textit{Mock examination working space.}\vfill
\clearpage\textbf{Draft Paper}\par\textit{Mock examination working space.}\vfill
\end{document}
